% GENERATED, edit the French unit or the translation table
\mxnewpage\section*{Solutions EV3-B}\label{EV3-B-CORR}
\textbf{MCQ:} 1 A, 2 B, 3 A, 4 B, 5 D, 6 B, 7 C, 8 C, 9 C, 10 D, 11 D, 12 D. Each correct answer: 0.5 point; total 6.
\paragraph{Checks of the fluency questions.} Order of operations; sum of fractions; subtraction of exponents; multiplying coefficient 1.2; division by 0.8; product 1.25 times 0.8; solving an affine equation; vertical increment divided by horizontal; mean; angle divided by 360; complement to 1; substitution in brackets.
\subsection*{Exercise 1 -- Four trials and a gain}
\begin{enumerate}
\item $(\dfrac{1}{2})^4=\dfrac{1}{16}$. \textit{Marking: 1 point. Result and basic justification expected; half a point possible for a correct but unfinished approach.}
\item $E(G)=\dfrac{1}{2}$, $E(G^2)=4$; $V(G)=\dfrac{15}{4}$. \textit{Marking: 2 points. Method or justification: 1; calculation and conclusion: 1.}
\item Stake $\dfrac{1}{2}$. The expectation of the gross gain is $\dfrac{1}{2}$, not possible in a single game. \textit{Marking: 2 points. Method or justification: 1; calculation and conclusion: 1.}
\end{enumerate}
\subsection*{Exercise 2 -- A derivative of constant sign}
\begin{enumerate}
\item $f\prime=(x^2+2x+1)e^x=(x+1)^2e^x$. \textit{Marking: 1 point. Result and basic justification expected; half a point possible for a correct but unfinished approach.}
\item Positive except at -1; f is strictly increasing (on both sides and by continuity). No extremum at -1. \textit{Marking: 2 points. Method or justification: 1; calculation and conclusion: 1.}
\item $f(0)=1,f\prime(0)=1$: $y=x+1$; $f(1)=2e$. \textit{Marking: 2 points. Method or justification: 1; calculation and conclusion: 1.}
\end{enumerate}
\subsection*{Exercise 3 -- Projection onto a diagonal}
\begin{enumerate}
\item Normal (1;-1). \textit{Marking: 1 point. Result and basic justification expected; half a point possible for a correct but unfinished approach.}
\item H(1;1), PH=$\sqrt2$. \textit{Marking: 1 point. Result and basic justification expected; half a point possible for a correct but unfinished approach.}
\item Angle $\dfrac{\pi}{4}$; vector $(\dfrac{\sqrt 2}{2};\dfrac{\sqrt 2}{2})$. \textit{Marking: 2 points. Method or justification: 1; calculation and conclusion: 1.}
\end{enumerate}
\paragraph{Calibration.} Target duration: 20--25 min for the fluency questions, 80--85 min for the exercises, 10--15 min for proofreading. These durations are editorial estimates to be measured in class. Accept any correct method; do not penalise twice a numerical error that is propagated consistently. The proposed marking scheme is not a decision of the team to harmonise marking.
